\section{Limitations}

We believe that the release of our multi-intent dataset framework, known as BlendX, represents a significant advancement in the field of multi-intent detection. However, we have identified several limitations in this work's current state and outline potential future directions to address these limitations.

A prominent issue in our work is our concentration on the MID problem, overlooking the task of slot filling. 
The challenge becomes particularly critical when utterances are modified after concatenation, emphasizing the need for more in-depth exploration in future research to accommodate the slot-filling task. 
Challenges seen in single-intent datasets, such as those in Banking77 with overlapping intents \cite{ying-thomas-2022-label}, continue to exist and highlight the need for more comprehensive dataset refinement strategies.
In Manual approach, we employed a variety of conjunctions; however, there is potential to improve creating more natural utterances by taking into account the relationships between sentences.
The metrics we have developed, focusing on utterance length, conjunctions, and pronouns, do not fully represent linguistic complexity, necessitating further enhancements. 

In conclusion, while our dataset marks a notable advancement in MID, several challenges persist. 
Addressing these will not only amplify our contributions but also serve the broader TOD field.